\chapter{L'apparato muscolo-scheletrico}

% ---------------------------------------------------------------------------
\section{Generalità sull'apparato muscolare}

Il termine \textbf{muscolo} deriva dal latino \textit{musculus} e indica un
organo composto in prevalenza da \textbf{tessuto muscolare}, ossia da
tessuto biologico dotato di capacità contrattile. I muscoli contribuiscono a
delimitare le regioni del corpo e sono gli effettori dei movimenti:
osservando un muscolo isolato si distinguono il \textbf{ventre} (la porzione
carnosa centrale), la \textbf{testa} e la \textbf{coda} (le due estremità),
raccordate all'osso tramite i \textbf{tendini}. Un muscolo scheletrico può
inoltre presentarsi in condizioni normali oppure atrofizzato, con evidente
riduzione del volume del ventre muscolare.

\subsection{Classificazione funzionale dei muscoli}

In base al ruolo svolto in un determinato movimento, i muscoli si
distinguono in:

\begin{itemize}
  \item \textbf{agonisti}: compiono il movimento;
  \item \textbf{antagonisti}: si oppongono al movimento prodotto dai muscoli
    agonisti;
  \item \textbf{sinergici}: concorrono allo stesso movimento.
\end{itemize}

La contrapposizione tra agonisti e antagonisti consente movimenti in
direzioni opposte e rende possibile il rallentamento di un'azione fino al
suo arresto. L'attività dei muscoli sinergici, invece, aumenta la forza di
contrazione, la velocità di movimento e ne modula l'intensità: si ottiene
così una gamma relativamente ampia di posizioni di un osso rispetto
all'altro.

% ---------------------------------------------------------------------------
\section{Muscoli della testa e del collo}

I muscoli della testa e del collo permettono i movimenti della faccia,
della lingua, della faringe e della laringe; sono coinvolti nella
comunicazione verbale e non verbale (ridere, parlare, accigliarsi,
sorridere, fischiettare) e nel nutrimento (masticazione, deglutizione). Si
dividono in due grandi gruppi: i muscoli che determinano i movimenti della
mandibola, detti \textbf{muscoli masticatori}, e i muscoli implicati nei
movimenti dei tegumenti facciali, detti \textbf{muscoli mimici} o
\textbf{pellicciai}.

Tra i principali muscoli della testa si riconoscono: il frontale, il
temporale, l'orbicolare dell'occhio, il buccinatore, l'orbicolare della
bocca, il digastrico, il tiroioideo, l'omoioideo, lo sternoioideo, gli
scaleni, il trapezio, l'elevatore della scapola, il massetere, il grande e
il piccolo zigomatico, il risorio, il quadrato del mento, il nasale, gli
elevatori del labbro superiore e delle labbra, e lo sternocleidomastoideo,
con i suoi due capi sternale e clavicolare.

\subsection{Muscoli masticatori}

Sono tutti muscoli pari. Il \textbf{temporale} si inserisce prossimalmente
sulla linea temporale inferiore e sulla squama dell'osso temporale, e
distalmente sul processo coronoideo della mandibola; innalza la mandibola e
chiude la bocca, retraendola e producendone i movimenti laterali. Il
\textbf{massetere} origina dall'arco zigomatico e si inserisce sull'angolo
mandibolare, con azione analoga di innalzamento della mandibola, chiusura
della bocca, protrusione, retrazione e movimenti laterali. Lo
\textbf{pterigoideo laterale} origina dalla grande ala dello sfenoide e dal
processo pterigoideo e si inserisce sul condilo mandibolare e
sull'articolazione temporomandibolare: la sua azione è l'apertura della
bocca, la protrusione della mandibola e i movimenti laterali. Lo
\textbf{pterigoideo mediale} origina dal processo pterigoideo, dall'osso
palatino e dalla tuberosità della mascella e si inserisce sull'angolo
mandibolare, innalzando la mandibola, chiudendo la bocca e producendo
movimenti laterali.

\subsection{Muscoli mimici}

I muscoli mimici sono innervati dal \textbf{nervo facciale}; la valutazione
delle loro funzioni comprende l'espressione facciale e il controllo delle
funzionalità proprie dei muscoli delle sopracciglia, delle palpebre, del
naso e della bocca. Il \textbf{temporo-parietale} origina dalla fascia
temporale sopra l'orecchio e si inserisce sulla galea aponeurotica. I
muscoli \textbf{estrinseci del padiglione auricolare} originano dalla galea
aponeurotica e si inseriscono sul padiglione stesso. Il \textbf{muscolo
nasale} presenta una parte trasversa, che origina dai gioghi alveolari dei
denti canini e incisivi superiori e si inserisce sul margine della narice,
e una parte alare, che origina dal margine dell'apertura piriforme della
mascella e si inserisce sul dorso e sull'ala del naso.

Tra gli altri muscoli mimici della faccia si riconoscono, con funzione di
origine e inserzione locale sui tegumenti: il ventre frontale
dell'occipito-frontale, il corrugatore del sopracciglio, l'orbicolare
dell'occhio, lo zigomatico minore e maggiore, l'elevatore del labbro
superiore, il procero, l'orbicolare della bocca, il risorio, il platisma,
il muscolo mentale, il depressore dell'angolo della bocca e il depressore
del labbro inferiore.

\subsection{Muscoli estrinseci dell'occhio}

I muscoli estrinseci dell'occhio governano i movimenti del bulbo oculare e
sono: il \textbf{retto inferiore} (azione principale: guardare in basso),
il \textbf{retto mediale} (rotazione mediale dell'occhio), il \textbf{retto
superiore} (guardare in alto), il \textbf{retto laterale} (rotazione
laterale dell'occhio), l'\textbf{obliquo inferiore} (rotazione oculare;
guardare in alto e lateralmente) e l'\textbf{obliquo superiore} (rotazione
oculare; guardare in basso e lateralmente). A questi muscoli si affianca
l'elevatore della palpebra superiore; i nervi coinvolti nella loro
innervazione sono l'oculomotore (III), il trocleare (IV), l'abducente (VI)
e il nervo ottico (II), che decorre insieme ai muscoli nella cavità
orbitaria.

\subsection{Muscoli della lingua}

I muscoli estrinseci della lingua sono: il \textbf{genioglosso}
(abbassamento e protrusione della lingua), lo \textbf{ioglosso}
(abbassamento e retrazione della lingua), il \textbf{palatoglosso}
(innalzamento della lingua, abbassamento del palato molle) e lo
\textbf{stiloglosso} (retrazione della lingua, innalzamento delle parti
laterali). Questi muscoli prendono rapporto con la mandibola, l'osso ioide
e il processo stiloideo dell'osso temporale.

\subsection{Muscoli anteriori del collo}

I \textbf{muscoli sopraioidei} comprendono: il \textbf{digastrico} (dal
greco \textit{dís}, doppio, e \textit{gastrós}, ventre), che origina con il
ventre anteriore dalla fossetta digastrica del margine inferiore del corpo
della mandibola e con il ventre posteriore dall'incisura mastoidea
dell'osso temporale, inserendosi sul corpo e sul grande corno dell'osso
ioide tramite un tendine intermedio: innalza l'osso ioide e abbassa la
mandibola, come nell'apertura della bocca; il ventre anteriore è innervato
da un ramo del nervo mandibolare (V paio di nervi cranici), il ventre
posteriore da un ramo del nervo faciale (VII paio); lo
\textbf{stiloioideo}, che origina dal processo stiloideo dell'osso
temporale e si inserisce sulla faccia anteriore del corpo dell'osso ioide,
innalzandolo e retraendolo, innervato da un ramo del nervo faciale; il
\textbf{miloioideo}, che dal greco \textit{mýle} (mola) origina dalla linea
miloioidea della faccia posteriore del corpo della mandibola e si inserisce
sulla faccia anteriore del corpo dell'osso ioide e sul rafe miloioideo,
innalzando l'osso ioide, sollevando il pavimento del cavo orale e la lingua
e abbassando la mandibola, innervato da un ramo del nervo mandibolare; il
\textbf{genioioideo}, che origina dall'apofisi genia inferiore della spina
mentale della mandibola e si inserisce sulla faccia anteriore del corpo
dell'osso ioide, innalzando l'osso ioide, tirandolo insieme alla lingua in
avanti e abbassando la mandibola, innervato dal primo paio di nervi
spinali cervicali (C1) per il tramite del nervo ipoglosso (XII paio).

I \textbf{muscoli sottoioidei} comprendono: l'\textbf{omoioideo}, che
origina dal margine superiore della scapola (medialmente all'incisura
scapolare) e dal legamento trasverso della scapola e si inserisce sul
margine inferiore del corpo dell'osso ioide e sul grande corno, abbassando
l'osso ioide e portandolo all'indietro, innervato dai rami anteriori dei
primi tre nervi cervicali del plesso cervicale (C1-C3); lo
\textbf{sternoioideo}, che origina dall'estremità mediale della clavicola e
dal manubrio dello sterno e si inserisce sul margine inferiore del corpo
dell'osso ioide, abbassandolo, con la medesima innervazione; lo
\textbf{sternotiroideo}, che origina dalla faccia posteriore del manubrio
dello sterno e dalla prima cartilagine costale e si inserisce sulla linea
obliqua della cartilagine tiroidea della laringe, abbassando la cartilagine
tiroidea e, tramite questa, l'osso ioide, con la medesima innervazione; il
\textbf{tiroioideo}, che origina dalla linea obliqua della cartilagine
tiroidea della laringe e si inserisce sul margine inferiore del corpo
dell'osso ioide e sul grande corno, innalzando la cartilagine tiroidea e
abbassando l'osso ioide, innervato dal primo paio di nervi spinali
cervicali (C1) per il tramite del nervo ipoglosso (XII paio).

\subsection{Muscolatura laterale del collo}

Lo \textbf{sternocleidomastoideo} è un muscolo bipennato, pari e
simmetrico, situato nella regione anterolaterale del collo, fusiforme e ben
visibile sotto la cute; possiede quattro capi ed è innervato dal nervo
accessorio. Origina dallo sterno e dalla clavicola (rispettivamente con un
capo sternale e un capo clavicolare che si fondono tra loro) e si inserisce
sul processo mastoideo e sulla linea nucale superiore. Compie un triplice
movimento: la rotazione della testa dalla parte opposta a quella della sua
contrazione, l'inclinazione dal lato della sua contrazione e l'estensione;
agisce inoltre come muscolo inspiratorio, sollevando lo sterno e la
clavicola.

Il \textbf{platisma} origina dalla seconda costa e si inserisce nella
regione masseterina, sulla mandibola e sulla commessura labiale; prende
origine dalla parte alta dei muscoli pettorali e dalla clavicola,
portandosi superiormente fino alla mandibola e alla commessura labiale: la
sua contrazione abbassa la mandibola e sposta in basso la commessura
labiale.

\subsection{Muscoli prevertebrali}

I muscoli prevertebrali flettono e inclinano lateralmente il capo. Il
\textbf{retto anteriore della testa} origina dall'atlante e si inserisce
sull'osso occipitale. Il \textbf{lungo del collo} origina dalle prime tre
vertebre toraciche e dalle ultime tre vertebre cervicali e si inserisce
sulle prime vertebre cervicali. Il \textbf{lungo della testa} origina
dall'osso occipitale e si inserisce sulle vertebre cervicali.

\subsection{Muscoli suboccipitali}

I muscoli suboccipitali agiscono come estensori del capo e rotatori; sono
quattro, pari e simmetrici, situati profondamente nella parte superiore del
rachide, in stretto contatto con il piano scheletrico. Collegano fra loro
le prime due vertebre cervicali e l'occipite, agendo sulle articolazioni
cranio-vertebrali: il \textbf{piccolo retto posteriore della testa}
collega il tubercolo posteriore dell'atlante all'osso occipitale, tra la
linea nucale superiore e quella inferiore; il \textbf{grande retto
posteriore della testa} collega il processo spinoso dell'epistrofeo
all'osso occipitale; l'\textbf{obliquo superiore della testa} si porta dal
processo trasverso dell'atlante all'osso occipitale, in prossimità del
processo mastoideo; l'\textbf{obliquo inferiore della testa} collega il
processo spinoso dell'epistrofeo al processo trasverso dell'atlante.

% ---------------------------------------------------------------------------
\section{Muscoli del torace}

I muscoli del torace si distinguono in \textbf{estrinseci}, che collegano
le ossa del torace a quelle della cintura pettorale, della cintura pelvica
e della colonna vertebrale, e \textbf{intrinseci}, che hanno origine e
inserzione nel torace stesso. Fanno parte dei muscoli estrinseci il gran
pettorale, il piccolo pettorale, il succlavio, il dentato anteriore e il
diaframma; fanno parte dei muscoli intrinseci gli intercostali interni,
esterni e intimi, il trasverso del torace, gli elevatori delle coste e i
sottocostali.

\subsection{Gran pettorale}

Il gran pettorale presenta tre parti: la parte \textbf{clavicolare}, che
origina dalla metà mediale del margine anteriore della clavicola; la parte
\textbf{sternocostale}, che origina dalla faccia anteriore dello sterno e
dalla seconda alla sesta cartilagine costale; la parte \textbf{addominale},
che origina dalla parte superiore della linea alba. I fasci si dirigono
lateralmente con un grosso tendine per inserirsi sulla cresta del tubercolo
maggiore dell'omero. Il muscolo adduce e ruota internamente l'omero; se
prende punto fisso sull'omero, solleva il tronco; agisce inoltre come
muscolo inspiratorio accessorio.

\subsection{Piccolo pettorale}

Origina dalla terza, quarta e quinta costa e si inserisce sul processo
coracoideo della scapola. Se prende punto fisso sul torace, abbassa la
scapola; se prende punto fisso sulla scapola, solleva le coste, agendo
quindi anche come muscolo inspiratorio.

\subsection{Succlavio}

È un piccolo muscolo triangolare che si estende fra la clavicola e la
prima costa, originando dalla prima costa e inserendosi sulla faccia
inferiore della clavicola. Può abbassare la spalla portando la clavicola
in basso e in avanti; se prende punto fisso sulla clavicola, innalza la
costa e contribuisce all'inspirazione. Fornisce inoltre protezione al
plesso brachiale e ai vasi sanguigni succlavi.

\subsection{Dentato anteriore}

Detto anche gran dentato o muscolo serrato anteriore, origina dalla faccia
laterale delle prime dieci coste e si inserisce sul margine mediale della
scapola. Determina lo spostamento in avanti della scapola; con punto fisso
sulla scapola, provoca l'elevazione delle coste, agendo come muscolo
inspiratorio.

\subsection{Diaframma}

Il diaframma è una lamina muscolo-fibrosa che separa la cavità toracica da
quella addominale, con la forma di una doppia cupola convessa
superiormente. Le sue origini, corrispondenti all'apertura inferiore del
torace, sono tripartite in una porzione sternale, una costale e una
lombare; si inserisce, a partire dal processo xifoideo, dalla settima alla
dodicesima costa e relativa cartilagine costale, e dalle vertebre lombari,
convergendo sul centro tendineo del diaframma. È il principale muscolo
respiratorio.

Nell'\textbf{inspirazione} l'escursione verso il basso del centro tendineo
del diaframma è favorita dal rilassamento della muscolatura addominale.
Raggiunto il massimo dell'escursione consentito dalla residua pressione
addominale, l'ulteriore contrazione delle fibre del diaframma, rispetto al
centro tendineo ormai bloccato, permette l'innalzamento delle coste, con
ampliamento dei diametri della gabbia toracica. L'\textbf{espirazione} si
avvale dell'azione dei muscoli intercostali interni e del diaframma; in
caso di espirazione forzata entrano in gioco anche i muscoli dell'addome
(retto dell'addome, obliquo esterno, obliquo interno e trasverso
dell'addome), che contraendosi determinano un aumento della pressione
intraddominale, sollevando ulteriormente il centro tendineo del diaframma e
abbassando le coste. Durante la fase inspiratoria che precede
l'espirazione forzata associata al soffiarsi il naso, al colpo di tosse o
allo starnuto, si verifica l'elevazione della gabbia toracica.

\subsection{Elevatori delle coste e sottocostali}

I muscoli \textbf{elevatori delle coste} hanno forma triangolare e sono
disposti lateralmente alla colonna; originano dai processi trasversi di C7
e di T9-T11 e si inseriscono sulla costa sottostante: sono muscoli
inspiratori, in quanto sollevano le coste. I muscoli \textbf{sottocostali}
originano dalla faccia interna delle coste e si inseriscono sulla faccia
interna della costa sottostante o della successiva: sono muscoli
espiratori, in quanto abbassano le coste.

\subsection{Muscoli intercostali}

I muscoli intercostali sono siti negli spazi intercostali e formano strati
sovrapposti che chiudono tali spazi dall'esterno all'interno. Gli
\textbf{intercostali esterni} originano dal margine inferiore delle coste e
si inseriscono sul margine superiore della costa sottostante; gli
\textbf{interni} originano dal labbro mediale del solco costale e si
inseriscono sul margine superiore della costa e della cartilagine
sottostante; gli \textbf{intimi} originano dal margine inferiore delle
coste e si inseriscono sul margine superiore della costa sottostante. La
contrazione degli intercostali esterni determina l'elevazione delle coste,
corrispondente all'inspirazione, ruotando l'asse lungo il collo delle
coste; la contrazione degli intercostali interni determina invece
l'abbassamento delle coste, corrispondente all'espirazione.

% ---------------------------------------------------------------------------
\section{Muscoli dell'addome}

I muscoli \textbf{anterolaterali} dell'addome svolgono un'importante
funzione nel determinare, con il loro tono, la pressione endoaddominale che
mantiene gli organi addominali in sede; inoltre, con la loro contrazione,
esercitano una compressione sugli stessi organi, facilitandone lo
svuotamento. Ne fanno parte il retto anteriore dell'addome, il piramidale,
l'obliquo interno, l'obliquo esterno e il trasverso dell'addome. I muscoli
\textbf{posteriori} comprendono invece il quadrato dei lombi e l'ileopsoas.

\subsection{Retto dell'addome}

È una spessa lamina muscolare che contribuisce a formare la parte mediana
anteriore della parete addominale, estendendosi dalla gabbia toracica alla
pelvi. Origina dal processo xifoideo e dalla quinta-settima cartilagine
costale e si inserisce sul ramo superiore del pube. Abbassa le coste
(azione espiratoria), flette il torace sulla pelvi o viceversa e aumenta la
pressione addominale. È il muscolo più famoso tra gli addominali, spesso
indicato con il termine informale \textit{«six pack»}.

\subsection{Obliquo esterno dell'addome}

La sua aponeurosi contribuisce, incontrando quella del muscolo omologo
dell'altro lato, a formare la \textbf{linea alba}. In alto si inserisce sul
processo xifoideo dello sterno; in basso, sulla cresta iliaca, sull'osso
iliaco e sul tubercolo pubico. Origina dalla faccia esterna e inferiore
delle ultime otto coste e si inserisce sulla cresta iliaca e
sull'aponeurosi. Abbassa le coste (azione espiratoria); i due muscoli
obliqui esterni possono agire in modo asimmetrico, determinando una
rotazione del tronco verso il lato opposto, oppure in sinergia, flettendo
il torace sulla pelvi.

\subsection{Obliquo interno dell'addome}

È posto sotto l'obliquo esterno e si allunga in direzione opposta a
quest'ultimo. Origina dal legamento inguinale, dalla spina iliaca
anteriore superiore, dalla cresta iliaca e dalla fascia toracolombare, e si
inserisce sulle ultime tre cartilagini costali e sull'aponeurosi. Ruota il
torace dal proprio lato e inclina la colonna; se si contrae da entrambi i
lati, flette la colonna, abbassa le coste (azione espiratoria) e aumenta
la pressione addominale.

\subsection{Piramidale}

È un piccolo muscolo posizionato nella parte inferiore del muscolo retto
dell'addome e situato nella sua guaina; è assente in un sesto dei
soggetti. Origina dal ramo superiore del pube e si inserisce sulla linea
alba: la sua contrazione tende la linea alba.

\subsection{Trasverso dell'addome}

È il più profondo tra i muscoli addominali, posto sotto i muscoli obliqui,
e avvolge la spina dorsale dandole protezione e stabilità; è un grande e
largo muscolo dell'addome. Origina dalle ultime sei cartilagini costali,
dalla fascia toracolombare, dalla cresta iliaca e dal legamento inguinale,
e si inserisce sull'aponeurosi. Abbassa le coste (azione espiratoria) e
aumenta la pressione addominale portando indietro le coste, favorendo lo
svuotamento dei visceri addominali e pelvici. Il corretto tempismo di
contrazione dei muscoli addominali (trasverso, obliquo interno, obliquo
esterno e retto dell'addome) è importante nella stabilizzazione della
colonna lombare e nella prevenzione della lombalgia.

\subsection{Quadrato dei lombi}

Origina dai processi costiformi di L2-L4 e dalla cresta iliaca, e si
inserisce sui processi costiformi di L1-L4 e sulla dodicesima costa.
Abbassa la dodicesima costa e inclina lateralmente la colonna lombare e la
pelvi.

\subsection{Grande psoas e iliaco}

Il \textbf{grande psoas} origina dai processi costiformi di L1-L5 e dai
corpi e dischi intervertebrali di T12 e di L1-L5, e si inserisce sul
piccolo trocantere del femore. Permette la flessione della coscia sul
bacino e la sua rotazione laterale; a femore fisso, flette il tronco e il
bacino. Il \textbf{muscolo iliaco} origina dalla fossa iliaca e si inserisce
anch'esso sul piccolo trocantere del femore: è un muscolo flessore che
completa l'azione del muscolo psoas; insieme, i due muscoli costituiscono
il complesso dell'\textbf{ileopsoas}.

% ---------------------------------------------------------------------------
\section{Muscoli del dorso}

La maggior parte del peso del corpo si sviluppa anteriormente alla colonna
vertebrale: per il suo sostegno (postura) e per il movimento, in
opposizione alla forza di gravità, sono necessari molti e potenti muscoli
che si inseriscono sui processi spinosi e trasversi delle vertebre. I
muscoli del dorso si suddividono in tre gruppi: gli \textbf{spino-appendicolari}
(trapezio, gran dorsale, piccolo e grande romboide, elevatore della
scapola), gli \textbf{spino-costali} (dentato posteriore superiore e
dentato posteriore inferiore) e gli \textbf{spino-dorsali}, a loro volta
articolati in spino-trasversari (tra cui lo splenio), erettore della
colonna vertebrale, trasverso-spinali, interspinosi e intertrasversari.

\subsection{Muscoli spino-appendicolari}

Il \textbf{trapezio} presenta tre parti. La parte \textbf{discendente}
origina dalla protuberanza occipitale esterna e dal legamento nucale e si
inserisce sulla clavicola. La parte \textbf{trasversa} origina dai processi
spinosi da C7 a T3 e si inserisce sull'acromion e sulla spina della
scapola. La parte \textbf{ascendente} origina dai processi trasversi da T3
a T12 e si inserisce sulla spina della scapola. Il trapezio eleva e adduce
la scapola oppure, se prende punto fisso sulla scapola, estende la testa
contraendosi bilateralmente, o la inclina dal proprio lato in contrazione
omolaterale.

Il \textbf{gran dorsale} origina dai processi spinosi di T6-T12 e da tutte
le vertebre lombari; i suoi fasci si portano in alto per inserirsi sul
tubercolo minore dell'omero. Adduce, estende e ruota internamente l'omero
e, se prende punto fisso sull'omero, eleva il tronco e le coste, come
avviene ad esempio nel movimento di arrampicata.

Il \textbf{piccolo romboide} origina dal legamento nucale e dai processi
spinosi di C7 e si inserisce sulla scapola; il \textbf{grande romboide}
origina dai processi spinosi di T1-T4 e si inserisce anch'esso sulla
scapola: entrambi spostano medialmente la scapola.

L'\textbf{elevatore della scapola} origina dai processi trasversi di
C1-C4, dirigendosi verso il basso e lateralmente per inserirsi sull'angolo
superiore della scapola: solleva la scapola e la sposta medialmente.

\subsection{Muscoli spino-costali}

I muscoli spino-costali occupano lo strato medio del dorso e sono
rappresentati dai muscoli dentati posteriori, superiore e inferiore. Il
\textbf{dentato posteriore superiore} origina dai processi spinosi di
C6-C7 e T1-T2 e si inserisce dalla seconda alla quinta costa, elevando le
coste; il \textbf{dentato posteriore inferiore} origina dai processi
spinosi di T11-T12 e L1-L2 e si inserisce sulle ultime quattro coste,
abbassandole. Entrambi partecipano ai movimenti respiratori.

\subsection{Muscoli spino-dorsali}

Tra gli spino-trasversari si distingue lo \textbf{splenio}: lo
\textbf{splenio della testa} origina dal legamento nucale e dai processi
spinosi di C7, T1 e T2, e si inserisce sul processo mastoideo del
temporale, inclinando la testa e ruotandola dal proprio lato, oppure, se
si contrae da entrambi i lati, estendendola; lo \textbf{splenio del
collo} estende la colonna cervicale.

Il gruppo dell'\textbf{erettore della colonna} (comprendente, tra gli
altri, il muscolo ileocostale e il muscolo lunghissimo nelle loro porzioni
cervicale, toracica e lombare) estende la colonna vertebrale, mantenendo
la stazione eretta; se si contrae da un solo lato, inclina la colonna.

I \textbf{trasverso-spinali} (comprendenti i muscoli semispinali della
testa, del collo e del torace, i muscoli rotatori lunghi e brevi e il
muscolo multifido) in contrazione bilaterale estendono e stabilizzano la
colonna vertebrale, mentre in contrazione unilaterale la inclinano dallo
stesso lato e la ruotano verso il lato opposto.

Gli \textbf{interspinosi} e gli \textbf{intertrasversari} sono sottili
fasci muscolari presenti nelle regioni cervicale, toracica e lombare (dove
sono più robusti) della colonna. I muscoli \textbf{interspinosi} agiscono
come estensori della colonna vertebrale; i muscoli \textbf{intertrasversari}
inclinano la colonna dal proprio lato e, se si contraggono
bilateralmente, la stabilizzano.

% ---------------------------------------------------------------------------
\section{Muscoli dell'arto superiore}

\subsection{Muscoli della spalla}

La spalla è un insieme funzionale di più articolazioni, che lavorano
insieme al fine di permettere la massima ampiezza dei movimenti omerali a
livello dell'articolazione scapolo-omerale: comprende nell'insieme cinque
articolazioni, di cui tre vere (diartrosi) e due spurie, cioè puramente
funzionali. I muscoli della spalla sono il deltoide, il sovraspinato,
l'infraspinato, il piccolo rotondo e il grande rotondo (questi ultimi
quattro, insieme al sottoscapolare, costituiscono la cuffia dei rotatori).

Il \textbf{deltoide} riveste la cuffia dei rotatori e origina con tre
capi: la parte clavicolare dalla porzione laterale del margine anteriore
della clavicola, la parte acromiale dall'apice e dal margine laterale
dell'acromion, la parte spinosa dal margine posteriore della spina della
scapola; si inserisce sulla tuberosità deltoidea dell'omero. Il fascio
anteriore produce flessione, adduzione e intrarotazione; il fascio
posteriore produce estensione, abduzione ed extrarotazione; il fascio
medio produce la massima abduzione.

La \textbf{cuffia dei rotatori} è un complesso muscolo-tendineo della
spalla che costituisce un importante mezzo di fissità e di
stabilizzazione dell'articolazione scapolo-omerale: i suoi grandi tendini
proteggono l'intera articolazione formando una vera e propria cuffia che
avvolge la testa dell'omero, coattandola nella cavità glenoidea. Il
\textbf{sottoscapolare} origina dalla fossa sottoscapolare e si inserisce
sul tubercolo minore dell'omero, adducendo e ruotando internamente
l'omero (intraruota/adduce). Il \textbf{sovraspinato} origina dai due
terzi mediali della fossa sopraspinata della faccia posteriore della
scapola e si inserisce sul tubercolo maggiore dell'omero, abducendo e
ruotando esternamente l'omero (abduce/extraruota). L'\textbf{infraspinato}
origina dalla fossa infraspinata della faccia posteriore della scapola e
si inserisce sul tubercolo maggiore dell'omero, ruotando esternamente
l'omero (extraruota). Il \textbf{piccolo rotondo} origina dalla metà
superiore della superficie ossea posta tra la fossa infraspinata e il
margine laterale della scapola e si inserisce sul tubercolo maggiore
dell'omero, ruotandolo esternamente, estendendolo e adducendolo
(extraruota/adduce). Il \textbf{grande rotondo} origina dall'angolo
inferiore della scapola e dalla metà inferiore della superficie ossea
posta tra la fossa infraspinata e il margine laterale della scapola, e si
inserisce sul tubercolo minore dell'omero: adduce l'omero, lo estende e
lo ruota internamente (adduce/intraruota).

\subsection{Muscoli anteriori del braccio}

I muscoli anteriori del braccio sono il bicipite brachiale, il
coracobrachiale e il brachiale. Il \textbf{bicipite brachiale} origina con
due capi: il capo lungo dal tubercolo sopraglenoideo della scapola, il
capo breve dall'apice del processo coracoideo della scapola; si inserisce
sulla tuberosità del radio. Sull'articolazione scapolo-omerale ha
funzione di stabilizzazione e partecipa alla flessione, all'abduzione e
all'adduzione del braccio; sull'articolazione del gomito flette
l'avambraccio sul braccio e, quando l'avambraccio è prono, agisce come
muscolo supinatore. Il \textbf{brachiale} origina dalla faccia
anteromediale e anterolaterale dell'omero, sotto la tuberosità deltoidea,
e si inserisce sulla tuberosità dell'ulna: flette l'avambraccio sul
braccio. Il \textbf{coracobrachiale} origina dal processo coracoideo
della scapola e si inserisce sulla faccia anteromediale del corpo
dell'omero: flette e adduce l'omero.

\subsection{Muscoli posteriori del braccio}

I muscoli posteriori del braccio sono il tricipite brachiale e l'anconeo.
Il \textbf{tricipite brachiale} origina con tre capi: il capo lungo dal
tubercolo sottoglenoideo della scapola, il capo laterale dalla faccia
posteriore dell'omero (sopra al solco del nervo radiale), il capo mediale
dalla faccia posteriore dell'omero (sotto al solco del nervo radiale); si
inserisce sull'olecrano dell'ulna. Sull'articolazione scapolo-omerale
adduce ed estende l'omero; sull'articolazione del gomito estende
l'avambraccio sul braccio. L'\textbf{anconeo} origina dall'epicondilo
laterale dell'omero e si inserisce sul margine laterale dell'olecrano
dell'ulna: estende l'avambraccio sul braccio.

\subsection{Cavità ascellare}

La cavità ascellare è uno spazio connettivale compreso tra il torace e la
parte superiore interna dell'arto superiore, di forma piramidale.
Presenta quattro pareti: la parete \textbf{anteriore}, formata dal
muscolo grande pettorale e, in basso, dai muscoli succlavio e piccolo
pettorale; la parete \textbf{posteriore}, formata dai muscoli
sottoscapolare, grande rotondo e grande dorsale; la parete
\textbf{laterale}, costituita dal capo breve del bicipite e dal muscolo
coracobrachiale; la parete \textbf{mediale}, costituita dal muscolo
dentato anteriore e dalla parete toracica. Presenta inoltre una
\textbf{base}, corrispondente alla fascia ascellare fissata ai margini
inferiori dei muscoli grande pettorale e grande dorsale (all'interno
della cavità, immersi nel tessuto adiposo, si trovano numerosi linfonodi
ascellari), e un \textbf{apice}, corrispondente allo spazio compreso tra
la clavicola, la prima costa e la base del processo coracoideo, a livello
del quale sono presenti i vasi succlavi che trapassano negli ascellari e
i rami del plesso brachiale.

\subsection{Muscoli dell'avambraccio}

I muscoli dell'avambraccio si distinguono in tre gruppi: quelli
\textbf{anteriori}, flessori e pronatori, detti epitrocleari; quelli
\textbf{laterali}, detti anche epitrocondiloidei; quelli
\textbf{posteriori}, estensori e supinatori, detti epicondiloidei.

I muscoli anteriori dell'avambraccio sono disposti su quattro strati. Il
\textbf{primo strato} comprende il pronatore rotondo, il flessore radiale
del carpo, il palmare lungo e il flessore ulnare del carpo; il
\textbf{secondo strato} comprende il flessore superficiale delle dita; il
\textbf{terzo strato} comprende il flessore profondo delle dita e il
flessore lungo del pollice; il \textbf{quarto strato} comprende il
pronatore quadrato.

Il \textbf{pronatore rotondo} origina con due capi, il capo omerale
dall'epicondilo mediale dell'omero e il capo ulnare dal processo
coronoideo dell'ulna, e si inserisce sulla faccia laterale del radio:
prona l'avambraccio ed è flessore dell'avambraccio sul braccio. Il
\textbf{flessore radiale del carpo} origina con due capi dall'epicondilo
mediale dell'omero e si inserisce con un lungo tendine alla base del
secondo osso metacarpale: flette e adduce la mano, flette l'avambraccio e
partecipa alla pronazione. Il \textbf{palmare lungo} origina
dall'epicondilo mediale dell'omero e si inserisce sull'aponeurosi
palmare della mano: flette la mano sull'avambraccio. Il \textbf{flessore
ulnare del carpo} origina con due capi, il capo omerale dall'epicondilo
mediale dell'omero e il capo ulnare dall'olecrano dell'ulna, e si
inserisce sull'osso pisiforme della mano: flette la mano sull'avambraccio
e la adduce.

Il \textbf{flessore superficiale delle dita} origina con due capi, il
capo omeroulnare dall'epicondilo mediale dell'omero e dal processo
coronoideo dell'ulna, il capo radiale dalla faccia anteriore prossimale
del radio, e si inserisce sulla seconda falange delle dita: flette le
falangi medie del secondo, terzo, quarto e quinto dito.

Il \textbf{flessore profondo delle dita} origina dalla porzione superiore
delle facce anteriore e mediale dell'ulna, dalla membrana interossea e
dal margine interosseo del radio, e si inserisce sulla faccia palmare
della falange distale: flette la falange distale del secondo, terzo,
quarto e quinto dito. Il \textbf{flessore lungo del pollice} origina
dalla faccia anteriore del radio, dalla membrana interossea e dal
processo coronoideo dell'ulna, e si inserisce sulla faccia palmare della
falange distale del pollice: ne flette la falange distale.

Il \textbf{pronatore quadrato} origina dal margine anteriore dell'ulna e
si inserisce sul margine anteriore del radio: ruota internamente
l'avambraccio nel movimento di pronazione.

I muscoli posteriori dell'avambraccio si articolano in un piano
superficiale e in un piano profondo. Nel \textbf{piano superficiale} si
trovano l'estensore delle dita, l'estensore del mignolo e l'estensore
ulnare del carpo. L'\textbf{estensore delle dita} origina dalla faccia
posteriore dell'epicondilo laterale dell'omero e raggiunge, con quattro
tendini finali, le inserzioni dorsali della seconda e terza falange del
secondo, terzo, quarto e quinto dito: estende le ultime quattro dita.
L'\textbf{estensore del mignolo} origina dall'epicondilo dell'omero e si
inserisce nell'aponeurosi della seconda e della terza falange del quinto
dito, fondendosi con l'aponeurosi dell'estensore comune delle dita:
estende il mignolo indipendentemente dalle altre dita. L'\textbf{estensore
ulnare del carpo} origina dalla faccia posteriore dell'epicondilo
dell'omero e si inserisce sulla faccia dorsale della base del quinto osso
metacarpale: estende e abduce la mano.

Nel \textbf{piano profondo} si trovano il supinatore, l'abduttore lungo
del pollice, l'estensore breve del pollice, l'estensore lungo del pollice
e l'estensore dell'indice. Il \textbf{supinatore} origina dall'epicondilo
dell'omero e si inserisce sul corpo del radio: porta l'avambraccio in
supinazione completa. L'\textbf{abduttore lungo del pollice} origina
dall'estremità superiore del radio e dell'ulna e dalla membrana
interossea, e si inserisce sulla base del primo osso metacarpale:
abduce il pollice. L'\textbf{estensore breve e lungo del pollice}
originano dal corpo dell'ulna e dalla membrana interossea
dell'avambraccio: il breve raggiunge ed estende la prima falange, il
lungo raggiunge ed estende la seconda falange del pollice. L'\textbf{estensore
dell'indice} origina dalla membrana interossea e il suo tendine si fonde
con il tendine dell'estensore comune delle dita, raggiungendo il secondo
dito.

I muscoli laterali dell'avambraccio sono il brachioradiale, l'estensore
radiale lungo del carpo e l'estensore radiale breve del carpo. Il
\textbf{brachioradiale} origina dal corpo dell'omero e si inserisce sul
processo stiloideo del radio: flette l'avambraccio sul braccio,
assumendo una posizione intermedia tra la pronazione e la supinazione.
L'\textbf{estensore radiale lungo del carpo} origina sopra l'epicondilo
dell'omero e si inserisce sul secondo osso metacarpale: estende e abduce
la mano. L'\textbf{estensore radiale breve del carpo} origina sopra
l'epicondilo dell'omero e si inserisce sul terzo osso metacarpale:
estende e abduce la mano.

\subsection{Muscoli della mano}

I muscoli della mano sono tutti situati nella faccia palmare e si
distinguono in tre gruppi: i muscoli dell'\textbf{eminenza tenar}
(laterali), i muscoli dell'\textbf{eminenza ipotenar} (mediali) e i
\textbf{muscoli palmari} (intermedi).

Fanno parte dei muscoli dell'eminenza tenar l'abduttore breve del
pollice, il flessore breve del pollice, l'opponente del pollice e
l'adduttore del pollice, che si inseriscono prossimalmente sul
retinacolo dei muscoli flessori e sui tubercoli delle ossa scafoide e
trapezio (o, nel caso dell'adduttore, sulle ossa trapezoide, capitato e
uncinato e sul margine anteriore del terzo osso metacarpale), e
distalmente sulla base della falange prossimale del pollice o sul primo
osso metacarpale.

Fanno parte dei muscoli dell'eminenza ipotenar il palmare breve,
l'abduttore del mignolo, il flessore breve del mignolo e l'opponente del
mignolo, che si inseriscono prossimalmente sull'aponeurosi palmare,
sull'osso pisiforme, sul legamento pisouncinato e sul tendine del
muscolo flessore ulnare del carpo, sul retinacolo dei muscoli flessori e
sull'osso uncinato, e distalmente sul derma dell'eminenza ipotenar, sulla
prima falange del mignolo o sul quinto osso metacarpale.

Fanno parte dei muscoli palmari i lombricali, che si inseriscono
prossimalmente sui tendini del muscolo flessore profondo delle dita e
distalmente sul rispettivo tendine del muscolo estensore delle dita; gli
interossei palmari, che si inseriscono prossimalmente sulla faccia
mediale del secondo osso metacarpale (il primo) o sulla faccia laterale
del quarto e del quinto osso metacarpale (il secondo e il terzo) e
distalmente sui tendini del muscolo estensore delle dita; e gli
interossei dorsali, che si inseriscono prossimalmente sulle ossa
metacarpali e distalmente sulla base della prima falange e sul tendine
del muscolo estensore delle dita.

% ---------------------------------------------------------------------------
\section{Muscoli dell'arto inferiore}

L'arto inferiore è l'appendice che si articola con il tronco nella sua
parte inferiore e serve alla locomozione. Si suddivide in sei regioni,
collegate da tre articolazioni principali. L'\textbf{anca} salda l'arto
inferiore al tronco per mezzo dell'articolazione coxo-femorale, tra osso
dell'anca e femore. La \textbf{coscia} è il segmento prossimale, lungo e
voluminoso, dotato di una possente muscolatura. Il \textbf{ginocchio}
contiene l'articolazione che connette il femore alla tibia, con la rotula
che ne stabilizza i movimenti, e comprende la regione poplitea. La
\textbf{gamba}, propriamente detta in termini anatomici, è il segmento
intermedio dell'arto inferiore e comprende una muscolatura particolarmente
sviluppata, posteriormente nel polpaccio e anteriormente nello stinco. La
\textbf{caviglia} comprende l'articolazione tibio-tarsica, o talo-crurale;
in superficie se ne osservano i rilievi dei malleoli mediale e laterale,
mentre posteriormente decorre il tendine calcaneale, o tendine di
Achille. Il \textbf{piede} è il terzo segmento, distale, e comprende le
dita e le articolazioni fra le ventisei ossa maggiori che lo strutturano,
insieme a un numero variabile di ossa sesamoidi e accessorie.

I muscoli dell'anca si distinguono in interni, rappresentati dall'ileo-psoas
e dal piccolo psoas, ed esterni, rappresentati dal grande, medio e piccolo
gluteo, dai gemelli, dall'otturatore esterno, dal piriforme, dal quadrato
del femore e dal tensore della fascia lata.

\subsection{Muscoli dell'anca}

Il \textbf{muscolo grande psoas} origina dalle superfici laterali della
dodicesima vertebra toracica e dalle prime quattro vertebre lombari; a
livello del bacino si congiunge con il \textbf{muscolo iliaco}, che origina
lateralmente dalla fossa iliaca, formando con esso un unico muscolo,
l'\textbf{ileo-psoas}, che si inserisce sul piccolo trocantere. Se prende
punto fisso sulla colonna vertebrale, flette la coscia sul bacino,
adducendola e intraruotandola; se prende punto fisso sul femore, flette il
tronco sul bacino e lo inclina dal proprio lato.

Il \textbf{piccolo psoas} origina dall'ultima vertebra toracica, dalla
prima vertebra lombare e dal disco interposto, e si inserisce
sull'eminenza ileopubica: partecipa alla flessione del tronco sul bacino.

Il \textbf{tensore della fascia lata} è un muscolo appiattito e allungato
che origina dalla spina iliaca anteriore superiore e dall'estremità
anteriore della cresta iliaca, e si inserisce sul condilo laterale della
tibia. Tende la fascia lata, flette e abduce la coscia contribuendo al
mantenimento della stazione eretta, ed estende inoltre la gamba sulla
coscia. Forma un fascio longitudinale che è muscolare solo nel primo
tratto, mentre nell'ultimo tratto diventa tendineo, costituendo il tratto
ileotibiale: la sua funzione è tendere la fascia lata, la lamina
connettivale che avvolge tutti i muscoli della coscia, e correggere la
posizione del ginocchio impedendo alla tibia di spostarsi medialmente.

Il \textbf{grande gluteo} è un muscolo quadrilatero posto superficialmente
nella regione glutea; origina dalla cresta iliaca, dalla linea glutea
posteriore e dalle superfici laterali del sacro e del coccige, e si
inserisce sulla tuberosità glutea del femore. Estende e ruota
esternamente la coscia, partecipando al mantenimento della stazione
eretta e alla deambulazione.

Il \textbf{medio gluteo} origina dalla cresta iliaca, dal tratto dell'osso
dell'anca compreso tra le linee glutee anteriore e posteriore e dalla
spina iliaca anteriore superiore, e si inserisce sul lato superomediale
del grande trocantere. Porta in abduzione la coscia ed è anche
extrarotatore; se prende punto fisso sul femore, inclina il bacino
lateralmente, e la sua contrazione bilaterale contribuisce al mantenimento
della stazione eretta.

Il \textbf{piccolo gluteo} ha una forma triangolare ed è posto in
profondità rispetto al medio gluteo; origina dalla linea glutea inferiore
e si inserisce sul grande trocantere, agendo come intrarotatore e
abduttore.

Il \textbf{piriforme} ha una forma triangolare ed è posto tra la
superficie esterna del sacro e il femore; origina dal sacro e dal margine
della grande incisura ischiatica e si inserisce sul grande trocantere.
Abduce la coscia ruotandola esternamente e stabilizza l'articolazione
coxo-femorale.

L'\textbf{otturatore esterno}, posto profondamente in prossimità
dell'articolazione dell'anca, origina dalla fossa otturatoria esterna e si
inserisce sul grande trocantere: è un muscolo stabilizzatore che ruota
esternamente la coscia. L'\textbf{otturatore interno}, posto all'interno
del bacino tra i due muscoli gemelli, origina dalla fossa otturatoria
interna e si inserisce anch'esso sul grande trocantere, agendo come
stabilizzatore e ruotando internamente la coscia.

I \textbf{gemelli}, superiore e inferiore, hanno direzione orizzontale;
il gemello superiore origina dal margine inferiore della spina ischiatica
e il gemello inferiore dalla tuberosità ischiatica, entrambi inserendosi
sulla fossetta trocanterica. Ruotano esternamente la coscia e
stabilizzano l'articolazione dell'anca.

Il \textbf{quadrato del femore} origina dalla tuberosità ischiatica e si
inserisce sulla cresta intertrocanterica: ruota la coscia in fuori ed è
abduttore.

\subsection{Muscoli della coscia}

I muscoli della coscia si distinguono in tre gruppi: gli \textbf{anteriori},
estensori, rappresentati dal quadricipite femorale e dal sartorio; i
\textbf{mediali}, adduttori, rappresentati dal gracile, dal pettineo,
dall'adduttore lungo, dall'adduttore breve e dal grande adduttore; i
\textbf{posteriori}, flessori, rappresentati dal bicipite femorale, dal
semitendinoso e dal semimembranoso.

Il \textbf{quadricipite femorale} origina con quattro capi. Il
\textbf{retto femorale} ha un decorso rettilineo e origina dalla spina
iliaca anteriore inferiore e dal contorno superiore dell'acetabolo. Il
\textbf{vasto laterale} origina dal grande trocantere, dalla tuberosità
glutea e dal labbro laterale della linea aspra. Il \textbf{vasto mediale}
origina dal collo anatomico del femore e dal labbro mediale della linea
aspra. Il \textbf{vasto intermedio} origina dalla superficie
antero-laterale del femore, arrivando anch'esso alla linea aspra. Tutti e
quattro i capi si inseriscono sulla patella.

Il \textbf{sartorio} è il muscolo più superficiale della coscia e ha un
decorso lateromediale e dall'avanti all'indietro, incrociando la coscia:
parte dall'anca in posizione laterale e arriva in posizione mediale al
ginocchio. Origina dalla spina iliaca anteriore superiore e si inserisce
sul condilo mediale della tibia, prendendo rapporti anche con il condilo
mediale del femore. Flette la gamba sulla coscia, flette la coscia sul
bacino ed extraruota la coscia abducendola.

Il \textbf{bicipite femorale} origina con due capi: il capo lungo dalla
tuberosità ischiatica, con un tendine comune al semitendinoso, e il capo
breve dalla linea aspra del femore; si inserisce sul condilo laterale
della tibia e sulla testa della fibula. Flette la gamba sulla coscia e la
ruota esternamente.

Il \textbf{semitendinoso} e il \textbf{semimembranoso} originano entrambi
dalla tuberosità ischiatica e si inseriscono entrambi sul condilo mediale
della tibia: estendono la coscia sul bacino, flettono la gamba sulla
coscia e la ruotano internamente.

Il \textbf{gracile} origina dal ramo ischio-pubico e si inserisce sul
condilo mediale della tibia: adduce la coscia, flette la gamba sulla
coscia e la ruota internamente.

Il \textbf{pettineo} origina dalla cresta pettinea dell'osso dell'anca e
si inserisce sulla linea pettinea del femore, agendo come rotatore
esterno, flessore e adduttore della coscia.

L'\textbf{adduttore lungo} origina dal ramo superiore del pube, tra la
sinfisi e il tubercolo pubico, e si inserisce sulla linea aspra del
femore. L'\textbf{adduttore breve} origina anch'esso dal ramo superiore
del pube, sotto l'adduttore lungo, e si inserisce parimenti sulla linea
aspra del femore. Entrambi agiscono come rotatori esterni, flessori e
adduttori della coscia.

Il \textbf{grande adduttore} origina dalla faccia anteriore della
tuberosità ischiatica; i suoi fasci superiori e medi si inseriscono sulla
linea aspra del femore, mentre i fasci inferiori si inseriscono al
tubercolo del grande adduttore. È rotatore esterno e adduttore della
coscia.

Un'immagine d'insieme dei muscoli dell'anca, della coscia e della gamba,
nelle visioni anteriore, posteriore e laterale dell'arto destro, mostra
inoltre, a completamento del quadro topografico, i muscoli propri della
gamba e del piede che saranno trattati separatamente: sulla faccia
anteriore il tibiale anteriore, l'estensore lungo delle dita e
l'estensore lungo dell'alluce; sulla faccia laterale i peronieri lungo,
breve e terzo; posteriormente il gastrocnemio con i capi mediale e
laterale, il soleo, il plantare, il flessore lungo delle dita e il
flessore lungo dell'alluce, con il tendine calcaneale e i retinacoli
superiore e inferiore dei muscoli estensori e dei muscoli peronieri.

\subsection{Triangolo femorale}

Il triangolo femorale è una regione di forma triangolare posta nella
porzione prossimale della faccia anteriore della coscia. È delimitato in
alto dal legamento inguinale, in basso e lateralmente dal muscolo
sartorio, in basso e medialmente dal muscolo adduttore lungo.
Superficialmente è rivestito dalla fascia lata, che centralmente prende
il nome di fascia cribrosa ed è attraversata da numerosi vasi, tra cui la
grande vena safena; il pavimento è delimitato lateralmente dal muscolo
ileopsoas e medialmente dal muscolo pettineo.

Le ernie crurali rappresentano la principale patologia di questa regione,
che offre inoltre un facile accesso alle arterie maggiori: l'angioplastica
coronarica e l'angioplastica periferica vengono spesso eseguite entrando
nell'arteria femorale. Il triangolo femorale contiene infine i bacini
superficiali e profondi dei linfonodi inguinali ed è la posizione ideale
per la linfoadenectomia inguinale.

\subsection{Fossa poplitea}

La fossa poplitea è una fossa situata posteriormente al ginocchio ed è
occupata da tessuto adiposo, nel quale sono contenuti i linfonodi
poplitei, il nervo ischiatico, la vena poplitea e l'arteria poplitea. Ha
forma di rombo: i suoi limiti sono, in alto e lateralmente, il muscolo
bicipite femorale, in alto e medialmente i tendini dei muscoli
semitendinoso e semimembranoso, in basso e medialmente il capo mediale del
muscolo gastrocnemio, in basso e lateralmente il capo laterale del
muscolo gastrocnemio.

Tra le patologie che possono interessare questa regione si ricordano le
cisti poplitee, gli aneurismi dell'arteria poplitea, gli aneurismi della
vena poplitea e la sindrome da intrappolamento dell'arteria poplitea.

\subsection{Muscoli della gamba}

I muscoli della gamba si distinguono in tre gruppi: gli \textbf{anteriori},
estensori, rappresentati dal tibiale anteriore, dall'estensore lungo delle
dita, dall'estensore lungo dell'alluce e dal peroniero anteriore; i
\textbf{laterali}, rappresentati dal peroniero lungo e dal peroniero
breve; i \textbf{posteriori}, flessori, rappresentati dal tricipite della
sura, dal plantare, dal popliteo, dal flessore lungo delle dita, dal
flessore lungo dell'alluce e dal tibiale posteriore.

Le strutture muscolari della gamba sono racchiuse nella fascia crurale,
continuazione distale della fascia lata; in prossimità della caviglia si
trovano i retinacoli che, assicurando i tendini in profondità,
ottimizzano la funzione della muscolatura crurale. A muovere la gamba è
principalmente la muscolatura della coscia: l'unica eccezione è
rappresentata dal muscolo popliteo, intrarotatore della gamba flessa,
situato nella loggia posteriore. Nel loro insieme, i muscoli della gamba
contribuiscono ai movimenti fondamentali della locomozione: la
plantarflessione, la dorsiflessione, l'estensione e la flessione della
gamba e delle dita dei piedi, l'eversione e l'inversione del piede.

Il \textbf{tibiale anteriore} origina dal condilo laterale della tibia e
dalla membrana interossea della gamba, e si inserisce sull'osso primo
cuneiforme e sulla base del primo metatarsale: flette dorsalmente il
piede e lo ruota medialmente.

L'\textbf{estensore lungo dell'alluce} origina dal corpo della fibula e
dalla membrana interossea, e si inserisce sulla faccia dorsale della
seconda falange dell'alluce: estende l'alluce e ruota medialmente il
piede adducendolo.

L'\textbf{estensore lungo delle dita} origina dal condilo laterale della
tibia, dalla membrana interossea e dalla testa della fibula; il suo
tendine si scompone in quattro tendini che raggiungono dorsalmente le
ultime quattro falangi: estende dorsalmente le dita e ruota in fuori il
piede.

Il \textbf{peroniero anteriore} origina dalla parte media del corpo della
fibula e dalla membrana interossea, e si inserisce sulla base del quinto
metatarsale: flette dorsalmente il piede ruotandolo in fuori.

Il \textbf{peroniero lungo} origina dalla faccia laterale della testa
della fibula e si porta sulla pianta del piede, inserendosi alla base del
primo osso metatarsale: flette plantarmente il piede, lo abduce e lo
ruota lateralmente. Il \textbf{peroniero breve} origina profondamente
rispetto al peroniero lungo e si inserisce sulla base del quinto osso
metatarsale: abduce il piede e lo ruota lateralmente.

Il \textbf{tricipite della sura} è una voluminosa massa carnosa che
determina la sporgenza del polpaccio; consta di tre capi, di cui due
superficiali, che formano nell'insieme il \textbf{muscolo gastrocnemio}, e
uno più profondo, che prende il nome di \textbf{muscolo soleo}. Il
gastrocnemio origina dalla faccia posteriore del femore, sopra i due
condili, e si inserisce sulla faccia posteriore del calcagno; il soleo
origina dalla tibia e dalla fibula e si inserisce sul tendine calcaneale.
Entrambi flettono plantarmente il piede e sollevano il corpo sulla punta
dei piedi.

Il \textbf{plantare} origina dal condilo laterale del femore e si
inserisce sulla tuberosità calcaneale insieme al tendine d'Achille:
partecipa all'azione del tricipite della sura.

Il \textbf{popliteo} è un muscolo corto e appiattito che si trova a
ridosso della faccia posteriore dell'articolazione del ginocchio: flette
la gamba e la ruota medialmente.

Il \textbf{flessore lungo delle dita} origina sotto la linea poplitea
della tibia; il suo tendine si scompone in quattro tendini che
raggiungono plantarmente le ultime quattro falangi: flette le falangi.

Il \textbf{flessore lungo dell'alluce} origina dalla membrana interossea
e dal corpo della fibula, e si porta sulla pianta del piede inserendosi
sulla faccia plantare della seconda falange dell'alluce: flette la
seconda falange dell'alluce.

Il \textbf{tibiale posteriore} origina dalla parte posteriore del corpo
della tibia, al di sotto della linea poplitea, e si inserisce sull'osso
navicolare e sul primo cuneiforme: flette plantarmente il piede e lo
adduce ruotandolo.
